\mode*
\section{Mehrdimensionale Funktionen}

\begin{frame}{Arten mehrdimensionaler Funktionen}
\begin{enumerate}
    \item Funktionen einer Variablen, die in $\R^n$ abbilden:
    \begin{gather*}
        \vf\colon \R\to \R^n,\qquad t\mapsto \vf(t)
    \end{gather*}
    \item Funktionen mehrerer Variablen, die nach $\R$ abbilden:
    \begin{gather*}
        f\colon \R^m\to \R,\qquad \vx \mapsto f(\vx)
    \end{gather*}
    \item Funktionen mehrerer Variablen, die in $\R^n$ abbilden:
    \begin{gather*}
        \vf\colon \R^m \to \R^n,\qquad \vx \mapsto \vf(\vx)
    \end{gather*}
\end{enumerate}
Notation: wir schreiben Vektorgrößen als $\vx,\vy,\vf$ zur Unterscheidung von Zahlen $t,x,f$.
\end{frame}

\subsection{Pfade und Wege}

%%%%%%%%%%%%%%%%%%%%%%%%%%%%%%%%%%%%%%%%%%%%%%%%%%%%%%%%%%%%
\begin{frame}
    \begin{Denkpause}
        Auf einer Wanderung bewegen Sie sich durch eine Landschaft. Was für ein mathematisches Objekt ist geeignet, diese Bewegung zu beschreiben.
    \end{Denkpause}
\end{frame}

%%%%%%%%%%%%%%%%%%%%%%%%%%%%%%%%%%%%%%%%%%%%%%%%%%%%%%%%%%%%
\begin{frame}{Analytische Eigenschaften von Pfaden}
    \begin{itemize}
        \item Stetigkeit, Differenzierbarkeit und Integrale von Pfaden sind komponentenweise definiert.
        \item Ableitungen
        \begin{gather*}
            \vf'(t) = \tfrac{d}{dt}
            \begin{pmatrix}
                \tfrac{d}{dt} f_1(t)\\
                \vdots\\
                \tfrac{d}{dt} f_n(t)
            \end{pmatrix}
            = \begin{pmatrix}
                f_1^(t)\\\vdots\\f_n'(t)
            \end{pmatrix}
        \end{gather*}
        \item Integrale
        \begin{gather*}
            \int \vf(t)\,dt = \vF(t),
        \qquad F_i(t) = \int f_i(t)\,dt,\quad i=1,\dots,n
        \end{gather*}
    \end{itemize}
\end{frame}

%%%%%%%%%%%%%%%%%%%%%%%%%%%%%%%%%%%%%%%%%%%%%%%%%%%%%%%%%%%%
\begin{frame}{Beispiele in zwei Dimensionen}
    \begin{columns}
        \begin{column}{.5\textwidth}
        \begin{block}{Geradlinige Bewegung}
            \begin{align*}
                x(t) &= f_1(t) = x_0 + t u\\
                y(t) &= f_2(t) = y_0+tv
            \end{align*}
            oder in Vektorschreibweise
            \begin{gather*}
                \vx(t) = \vx_0+t\vv
            \end{gather*}
            Ableitung
            \begin{gather*}
                \vx'(t) = \vv
            \end{gather*}
        \end{block}         
        \end{column}
        \pause
        \begin{column}{.5\textwidth}
        \begin{block}{Kreisbewegung}
            \begin{align*}
                x(t) &= r \cos t\\
                y(t) &= r \sin t
            \end{align*}
            mit der Ableitung
            \begin{align*}
                x(t) &= -r \sin t\\
                y(t) &= r \cos t
            \end{align*}
        \end{block}         
        \end{column}
    \end{columns}
\end{frame}

%%%%%%%%%%%%%%%%%%%%%%%%%%%%%%%%%%%%%%%%%%%%%%%%%%%%%%%%%%%%


\begin{envgknote}{hinzufügen}
   \begin{itemize}
       \item Wege: parametrische und nullstellen-Darstellung
       \item Begriffsklärung
       \item Beispiele Gerade und Kreis
       \item zu einem Weg kann es mehrere Pfade geben, die ihn parametrisieren
       \item Beispiel chemische Reaktion oder Biosystem
       \item Wann ist ein Weg als $y=f(x)$ darstellbar?
       \item Implizite Differentiation
   \end{itemize} 
\end{envgknote}
\mode<all>
